\section{Problem Formulation}
\label{sec:formulation}

\paragraph{Streaming setting.} Following the online evaluation protocol of
StreamBench \citep{wu2024streambench}, an agent faces an input--feedback sequence
$(x_1, fb_1), (x_2, fb_2), \dots, (x_T, fb_T)$ drawn from a task distribution and
presented in a fixed, seeded order. At step $t$ the agent produces
\begin{equation}
\hat y_t \;=\; f\!\big(p(x_t,\, r(M_t)) \,\big|\, \theta_t\big),
\end{equation}
where $\theta_t$ are the current model weights, $M_t$ is an external memory,
$r(\cdot)$ a retriever, and $p(\cdot)$ a prompting template. The environment
returns \emph{binary verifiable feedback}
$fb_t = g(x_t, \hat y_t) \in \{0, 1\}$---execution of generated code against unit
tests, execution-match of SQL, exact/label match for QA and diagnosis. Gold
references exist only inside $g$; the agent observes the bit. After feedback, the
agent may update any of $(p, r, M, \theta)$; the methods we study differ precisely
in \emph{which} of these they update (Table~\ref{tab:taxonomy}).

\paragraph{Evaluation.} Performance is \emph{prequential}: $\hat y_t$ is produced
and scored before $x_t$ influences any component, and the stream is consumed once.
We report final cumulative accuracy
$\mathrm{Acc}_T = \tfrac{1}{T}\sum_{t=1}^{T} h(\hat y_t, y_t)$
(each task's native metric $h$) and, equivalently, cumulative regret
$R_T = \sum_{t=1}^{T} (1 - h(\hat y_t, y_t))$, which penalizes late learning.
All methods on a given task see the identical stream order.

\paragraph{Why this setting is hard.} The agent gets one attempt per problem, a
single bit per attempt, no gold text ever, and no second pass over the stream.
Improvements must therefore be bootstrapped from the agent's \emph{own} verified
outputs---and any weight update risks destabilizing performance on the remainder
of the stream, a risk realized dramatically by the RL baseline in
\S\ref{sec:exp-fused}.

\begin{table}[t]
\centering
\caption{Method taxonomy: what each method updates and what signal it consumes.
\method{} is the only method that both updates weights and exploits a dense,
token-level training signal derived from verified self-generated text.}
\label{tab:taxonomy}
\small
\begin{tabular}{lcccc}
\toprule
Method & Updates $M$ & Updates $\theta$ & Training signal & Signal density \\
\midrule
Zero-shot (Base) & -- & -- & -- & -- \\
Self-StreamICL & \checkmark & -- & -- (retrieval only) & -- \\
REINFORCE++ & -- & \checkmark & scalar reward & per sequence \\
\textbf{\method{} (ours)} & \checkmark & \checkmark & verified self-output & per token \\
\bottomrule
\end{tabular}
\end{table}
